\documentclass[12pt,a4paper]{report}

% Russian language support
\usepackage{fontspec}
\usepackage{polyglossia}
\setdefaultlanguage{russian}
\setotherlanguage{english}

% Fonts with Cyrillic support
\setmainfont{DejaVu Serif}
\setsansfont{DejaVu Sans}
\setmonofont{DejaVu Sans Mono}

% Packages for enhanced document
\usepackage{geometry}
\usepackage{amsmath}
\usepackage{amsfonts}
\usepackage{amssymb}
\usepackage{graphicx}
\usepackage{fancyhdr}
\usepackage{titlesec}
\usepackage{tocloft}
\usepackage{xcolor}
\usepackage{listings}
\usepackage{tikz}
\usepackage{pgfplots}
\usepackage{float}

% Page setup
\geometry{margin=2.5cm}

% Header and footer
\pagestyle{fancy}
\fancyhf{}
\fancyhead[L]{Подготовка к экзамену}
\fancyhead[R]{\thepage}
\fancyfoot[C]{\textit{Учебные материалы}}

% Chapter styling
\titleformat{\chapter}[display]
  {\normalfont\huge\bfseries\color{blue!60!black}}
  {\chaptertitlename\ \thechapter}{20pt}{\Huge}

% Document metadata
\title{\textbf{Подготовка к экзамену}\\\large Конспект лекций и формулы}
\author{Студент}
\date{\today}

\begin{document}

% Title page
\begin{titlepage}
\centering
\vspace*{2cm}

{\huge\bfseries Подготовка к экзамену\par}
\vspace{1cm}
{\Large Конспект лекций и основные формулы\par}
\vspace{2cm}

{\large Предмет: \textbf{Математический анализ}\par}
\vspace{0.5cm}
{\large Курс: \textbf{1-й курс}\par}
\vspace{0.5cm}
{\large Семестр: \textbf{Весенний 2024}\par}

\vspace{3cm}

{\Large Студент: \textbf{Иван Петров}\par}
\vspace{0.5cm}
{\large Группа: \textbf{МТ-101}\par}
\vspace{0.5cm}
{\large Преподаватель: \textbf{Профессор Сидоров А.И.}\par}

\vfill

{\large Москва\par}
{\large \today\par}
\end{titlepage}

% Table of contents
\tableofcontents
\newpage

% Introduction
\chapter*{Введение}
\addcontentsline{toc}{chapter}{Введение}

Этот конспект содержит основные темы курса математического анализа, необходимые для подготовки к экзамену. Каждая глава соответствует отдельной экзаменационной теме.

\textbf{Структура документа:}
\begin{itemize}
    \item Теоретический материал
    \item Основные формулы и определения
    \item Примеры решения задач
    \item Графики и иллюстрации
\end{itemize}

% Chapter 1: Limits
\chapter{Пределы функций}

\section{Определение предела}
Предел функции $f(x)$ при $x \to a$ равен $L$, если:
\begin{equation}
\lim_{x \to a} f(x) = L \iff \forall \varepsilon > 0 \, \exists \delta > 0 : |x - a| < \delta \Rightarrow |f(x) - L| < \varepsilon
\end{equation}

\section{Основные пределы}
\begin{align}
\lim_{x \to 0} \frac{\sin x}{x} &= 1 \\
\lim_{x \to \infty} \left(1 + \frac{1}{x}\right)^x &= e \\
\lim_{x \to 0} \frac{e^x - 1}{x} &= 1 \\
\lim_{x \to 0} \frac{\ln(1 + x)}{x} &= 1
\end{align}

\section{Пример решения}
\textbf{Задача:} Найти $\lim_{x \to 0} \frac{\sin 3x}{2x}$

\textbf{Решение:}
\begin{align}
\lim_{x \to 0} \frac{\sin 3x}{2x} &= \lim_{x \to 0} \frac{3 \sin 3x}{6x} \cdot \frac{3}{2} \\
&= \frac{3}{2} \lim_{x \to 0} \frac{\sin 3x}{3x} \\
&= \frac{3}{2} \cdot 1 = \frac{3}{2}
\end{align}

% Chapter 2: Derivatives
\chapter{Производные функций}

\section{Определение производной}
Производная функции $f(x)$ в точке $x_0$:
\begin{equation}
f'(x_0) = \lim_{h \to 0} \frac{f(x_0 + h) - f(x_0)}{h}
\end{equation}

\section{Таблица производных}
\begin{center}
\begin{tabular}{|c|c|}
\hline
\textbf{Функция} & \textbf{Производная} \\
\hline
$x^n$ & $nx^{n-1}$ \\
$e^x$ & $e^x$ \\
$\ln x$ & $\frac{1}{x}$ \\
$\sin x$ & $\cos x$ \\
$\cos x$ & $-\sin x$ \\
$\tan x$ & $\sec^2 x$ \\
\hline
\end{tabular}
\end{center}

\section{Правила дифференцирования}
\begin{itemize}
    \item $(u + v)' = u' + v'$ — производная суммы
    \item $(uv)' = u'v + uv'$ — производная произведения
    \item $\left(\frac{u}{v}\right)' = \frac{u'v - uv'}{v^2}$ — производная частного
    \item $(f(g(x)))' = f'(g(x)) \cdot g'(x)$ — правило цепочки
\end{itemize}

% Chapter 3: Integrals  
\chapter{Интегралы}

\section{Неопределённый интеграл}
Неопределённый интеграл функции $f(x)$:
\begin{equation}
\int f(x) dx = F(x) + C
\end{equation}
где $F'(x) = f(x)$ и $C$ — константа интегрирования.

\section{Основные интегралы}
\begin{align}
\int x^n dx &= \frac{x^{n+1}}{n+1} + C \quad (n \neq -1) \\
\int \frac{1}{x} dx &= \ln|x| + C \\
\int e^x dx &= e^x + C \\
\int \sin x dx &= -\cos x + C \\
\int \cos x dx &= \sin x + C
\end{align}

\section{Определённый интеграл}
Формула Ньютона-Лейбница:
\begin{equation}
\int_a^b f(x) dx = F(b) - F(a)
\end{equation}

% Chapter 4: Series
\chapter{Ряды}

\section{Степенные ряды}
Степенной ряд имеет вид:
\begin{equation}
\sum_{n=0}^{\infty} a_n (x - c)^n = a_0 + a_1(x-c) + a_2(x-c)^2 + \ldots
\end{equation}

\section{Разложения в ряд Тейлора}
\begin{align}
e^x &= \sum_{n=0}^{\infty} \frac{x^n}{n!} = 1 + x + \frac{x^2}{2!} + \frac{x^3}{3!} + \ldots \\
\sin x &= \sum_{n=0}^{\infty} \frac{(-1)^n x^{2n+1}}{(2n+1)!} = x - \frac{x^3}{3!} + \frac{x^5}{5!} - \ldots \\
\cos x &= \sum_{n=0}^{\infty} \frac{(-1)^n x^{2n}}{(2n)!} = 1 - \frac{x^2}{2!} + \frac{x^4}{4!} - \ldots
\end{align}

% Add a simple graph using TikZ
\section{График функции}
\begin{figure}[H]
\centering
\begin{tikzpicture}
\begin{axis}[
    xlabel={$x$},
    ylabel={$y$},
    domain=-2*pi:2*pi,
    samples=100,
    grid=major,
    title={График функции $y = \sin x$}
]
\addplot[blue, thick] {sin(deg(x))};
\end{axis}
\end{tikzpicture}
\caption{Синусоида}
\end{figure}

% Conclusion
\chapter*{Заключение}
\addcontentsline{toc}{chapter}{Заключение}

Данный конспект охватывает основные темы курса математического анализа. Для успешной сдачи экзамена рекомендуется:

\begin{enumerate}
    \item Выучить все основные формулы
    \item Решить дополнительные задачи
    \item Повторить теоретический материал
    \item Практиковаться в вычислениях
\end{enumerate}

\textbf{Удачи на экзамене!}

\end{document}